\documentclass[11pt,a4paper]{article}

\usepackage[margin=2.5cm]{geometry}
\usepackage{amsmath,amssymb,amsthm}
\usepackage{bm}
\usepackage{booktabs}
\usepackage{xcolor}
\usepackage{hyperref}
\usepackage{enumitem}

\newtheorem{theorem}{Theorem}[section]
\newtheorem{lemma}[theorem]{Lemma}
\newtheorem{corollary}[theorem]{Corollary}
\newtheorem{proposition}[theorem]{Proposition}
\newtheorem{remark}[theorem]{Remark}
\newtheorem{assumption}[theorem]{Assumption}

\theoremstyle{definition}
\newtheorem{definition}[theorem]{Definition}

\newcommand{\R}{\mathbb{R}}
\newcommand{\dom}{\Omega}
\newcommand{\pore}[1]{\Omega_{#1}}
\newcommand{\iface}{\gamma}
\newcommand{\vel}{\bm{u}}
\newcommand{\pres}{p}
\newcommand{\stress}{\bm{\sigma}}
\newcommand{\normal}{\bm{n}}
\newcommand{\trac}{\bm{\lambda}}
\newcommand{\DtN}{\mathbf{G}}
\newcommand{\Schur}{\mathbf{S}}
\newcommand{\Shat}{\widehat{\Schur}}
\newcommand{\Shatuu}{\widehat{\Schur}{}^{uu}}
\newcommand{\Fhat}{\widehat{\bm{F}}}
\newcommand{\one}{\mathbf{1}}
\newcommand{\zero}{\bm{0}}
\newcommand{\Range}{\operatorname{Range}}
\newcommand{\Ker}{\operatorname{Ker}}
\newcommand{\rank}{\operatorname{rank}}
\newcommand{\modeset}{\mathcal{M}}
\newcommand{\loadvec}{\bm{f}}
\newcommand{\Vspace}{\mathcal{V}}
\newcommand{\Wspace}{\mathcal{W}}
\newcommand{\Lspace}{\Lambda}
\newcommand{\Pcal}{\mathcal{P}}

\hypersetup{colorlinks=true,linkcolor=blue,urlcolor=blue,citecolor=blue}

\begin{document}

\title{\textbf{Theory of the Affine-DDPNM}\\[4pt]
       \large Best Approximation, Nested Decomposition, and Error Separation\\[4pt]
       \large Version 3 --- corrected 2026-08-07}
\author{}
\date{}
\maketitle

\begin{center}
\fbox{\begin{minipage}{0.92\textwidth}
This document develops the Galerkin projection theory of the
Affine-DDPNM.  Starting from the broken-pressure
continuous-traction (C-DD) reference system, we prove:
(i)~the Schur solution at each modal level is the
$\Shatuu$-seminorm best approximation to the full C-DD trace solution;
(ii)~for nested trace subspaces the error admits an orthogonal
Pythagorean decomposition; and (iii)~the total volume error is bounded
in terms of a C-DD discretisation component and an
interface-modal-reduction component.  The reduction component is
identified exactly in the broken energy and broken-$H^1$ measures;
in $L^2$, an exact two-sided triangle bound is combined with a
one-sided Poincar\'e--Friedrichs estimate.

\textbf{Scope.}  The reference system throughout is the
broken-pressure C-DD formulation (discontinuous $P_1$ pressure across
pores).  Equivalence with the monolithic continuous-pressure
Taylor--Hood system is not established here; consequently the
discretisation term $\|\vel - \vel_h\|$ appearing in the separation
bounds is a C-DD discretisation error, not a standard FEM error.
All results assume planar interfaces with a fixed orthonormal frame.
\end{minipage}}
\end{center}

% ======================================================================
\section{C-DD reference system and trace subspaces}
% ======================================================================

\subsection{Notation and standing assumptions}

We work on a fixed finite-element mesh of the fluid domain
$\dom\subset\R^3$, partitioned into $N$ non-overlapping pore subdomains
$\{\pore{i}\}_{i=1}^N$ with $m$ planar internal interfaces
$\{\iface_k\}_{k=1}^m$.  On each pore we use the Taylor--Hood
$[P_2]^3$--$P_1$ pair with \textbf{discontinuous} pressure across
interfaces (the broken-pressure C-DD formulation).  All objects are
finite-dimensional; $h$ refers to the given mesh.

\begin{assumption}[Local algebraic stability]\label{ass:local}
  For every pore $i$:
  \begin{enumerate}[label=(\alph*),nosep,leftmargin=*]
    \item $A_i \in \R^{n_{u,i}\times n_{u,i}}$ is symmetric positive definite;
    \item $B_i \in \R^{n_{p,i}\times n_{u,i}}$ has full row rank;
    \item the constant function $1$ belongs to the local pressure space.
  \end{enumerate}
\end{assumption}

For the standard conforming velocity space, Assumption~\ref{ass:local}(a)
follows, for example, if $\pore{i}$ is connected and Lipschitz and its
Dirichlet/no-slip wall portion $\Gamma_{w,i}\subset\partial\pore{i}$ has
positive surface measure.  These geometric conditions exclude nonzero
constant-velocity null modes via the Poincar\'e--Friedrichs inequality.

\begin{assumption}[Matching interface trace spaces]\label{ass:matching}
  For every internal interface $\iface_k$ shared by pores $i$ and $j$,
  the two local velocity trace spaces coincide with a shared scalar
  nodal basis $\{\psi_{k,\ell}\}_{\ell=1}^{m_k}$ and a one-to-one DOF
  correspondence between the two sides.
\end{assumption}

\medskip\noindent
\textbf{Dimension conventions.}
$m_k$ = number of scalar trace DOFs on interface $k$;
$m_\Gamma = 3\sum_{k=1}^m m_k$ = total C-DD interface DOFs
(three components per trace node).
For pore $i$,
$r_i^u = 3\sum_{\gamma\in\mathcal{U}_i} m_{i,\gamma}$
(C-DD internal trace DOFs) and
$r_i^k = 3\sum_{\gamma\in\mathcal{K}_i} m_{i,\gamma}$
(C-DD boundary trace DOFs).

\subsection{Interface frames and affine scalar shapes}

On interface $\iface_k$ fix an orthonormal frame
$\{\normal_k,\bm{t}_{k,1},\bm{t}_{k,2}\}$ with $\normal_k$ directed
from the first pore of the pair to the second.  Let
$\bm{c}_k$ be the interface centroid and define the half-span-normalised
in-plane coordinates
\begin{equation}\label{eq:st}
  s_k(\bm{x}) = \frac{(\bm{x}-\bm{c}_k)\cdot\bm{t}_{k,1}}{a_k},\qquad
  t_k(\bm{x}) = \frac{(\bm{x}-\bm{c}_k)\cdot\bm{t}_{k,2}}{b_k},
\end{equation}
where $a_k,b_k > 0$ are the maximal absolute in-plane coordinate
values on $\iface_k$.  The three affine scalar shape functions are
\begin{equation}\label{eq:shapes}
  \varphi_0(\bm{x}) \equiv 1,\qquad
  \varphi_s(\bm{x}) = s_k(\bm{x}),\qquad
  \varphi_t(\bm{x}) = t_k(\bm{x}),
\end{equation}
and the mode set is
$\modeset := \{0,s,t\} \times \{n,t_1,t_2\}$.

\subsection{C-DD coupling, Schur complement, and energy seminorm}

For pore $i$ incident on $\iface_k$, define the oriented
inward-traction directions
$\bm{d}_{i,k,\beta}$ ($\beta\in\{n,t_1,t_2\}$) satisfying
$\bm{d}_{j,k,\beta} = -\bm{d}_{i,k,\beta}$ for the opposite pore $j$.

The C-DD coupling block for direction $\beta$ on $\iface_k$ is
\begin{equation}\label{eq:N}
  [N_{i,k}^{\beta}]_{\ell,j}
  = \int_{\iface_k} \psi_{k,\ell}\;
    \bm{d}_{i,k,\beta}\cdot\bm{\phi}_j^{(i)}\,\mathrm{d}S
  \;\in\; \R^{m_k\times n_{u,i}}.
\end{equation}
Collect the three-direction blocks on all internal ports of pore $i$
into $N_i^u \in \R^{r_i^u \times n_{u,i}}$, and on all boundary ports
into $N_i^k \in \R^{r_i^k \times n_{u,i}}$.  Let $\bm{\xi}^k \in
\R^{m_\Gamma^k}$ (with $m_\Gamma^k = \sum_i r_i^k$) denote the
globally prescribed boundary traction nodal values.  Through the
boundary affine restriction (see the companion
mathematical-properties document, Appendix~A) these are determined by
the prescribed inlet/outlet pressures.

The velocity response operator is
$H_i = A_i^{-1} - A_i^{-1}B_i^\top M_i^{-1}B_i A_i^{-1}$,
with $M_i = B_i A_i^{-1}B_i^\top$.
$H_i$ is symmetric positive semidefinite and satisfies
$H_i^\top A_i H_i = H_i$,
$\Ker(H_i) = \Range(B_i^\top)$.

The \textbf{internal-unknown block} of the local C-DD Schur complement is
\begin{equation}\label{eq:Shat-local}
  \Shatuu_i = N_i^u H_i (N_i^u)^\top \;\in\; \R^{r_i^u\times r_i^u}.
\end{equation}
Let $\widehat{R}_i^u : \R^{m_\Gamma} \to \R^{r_i^u}$ be the Boolean
restriction extracting the C-DD DOFs of pore $i$ from the global
vector of unknown interface DOFs.  The assembled C-DD Schur complement
(uu-block) is
\begin{equation}\label{eq:Shat}
  \Shatuu = \sum_{i=1}^N (\widehat{R}_i^u)^\top \Shatuu_i \widehat{R}_i^u
  \;\in\; \R^{m_\Gamma\times m_\Gamma}.
\end{equation}
$\Shatuu$ is symmetric positive semidefinite.

The C-DD right-hand side, with prescribed boundary tractions absorbed,
is
\begin{equation}\label{eq:Fhat}
  \Fhat = -\sum_{i=1}^N (\widehat{R}_i^u)^\top
    N_i^u H_i (N_i^k)^\top \widehat{R}_i^k \bm{\xi}^k
  \;\in\; \R^{m_\Gamma},
\end{equation}
where $\widehat{R}_i^k$ extracts the boundary DOFs of pore $i$ from
$\bm{\xi}^k$.

The \textbf{full C-DD system} for the unknown interface DOFs is
\begin{equation}\label{eq:full-sys}
  \Shatuu\,\bm{\xi} = \Fhat,\qquad \bm{\xi}\in\Vspace:=\R^{m_\Gamma}.
\end{equation}
\begin{assumption}[Full C-DD consistency]\label{ass:consistency}
  The assembled full C-DD Schur system~\eqref{eq:full-sys} is consistent,
  equivalently
  $\Fhat\in\Range(\Shatuu)$.
\end{assumption}
This assumption is the only global solvability hypothesis used in the
projection arguments below.  Its derivation from a fully specified
broken-pressure mixed Stokes formulation is not repeated in this
document.
Define the energy seminorm on $\Vspace$ by
$\|\bm{\xi}\|_{\Shatuu}^2 := \bm{\xi}^\top \Shatuu \bm{\xi}$.

\subsection{Restriction operators and nested trace subspaces}

\begin{definition}[Per-interface restriction matrices]\label{def:Pkl}
  For each interface $\iface_k$, let
  $[M_k]_{\ell,\ell'} = \int_{\iface_k}\psi_{k,\ell}\psi_{k,\ell'}\,\mathrm{d}S$
  be the face mass matrix and
  $[b_k^\alpha]_\ell = \int_{\iface_k} \varphi_\alpha\psi_{k,\ell}\,\mathrm{d}S$.
  Since $\{\psi_{k,\ell}\}_{\ell=1}^{m_k}$ is linearly independent in
  $L^2(\iface_k)$, $M_k$ is symmetric positive definite and hence
  invertible.  Define the $L^2$-projection vectors
  $p_k^\alpha = M_k^{-1} b_k^\alpha \in \R^{m_k}$ for
  $\alpha\in\{0,s,t\}$.
  For level $\ell\in\{0,1,2\}$ with $d_\ell$ coefficient DOFs per interface,
  \begin{itemize}
    \item $\ell=0$ (classic P0, $d_0=1$):
      $P_{k,0} = \begin{bmatrix} p_k^0 \\ \zero \\ \zero \end{bmatrix}
       \in \R^{3m_k\times 1}$;
    \item $\ell=1$ (constant-vector P0, $d_1=3$):
      $P_{k,1} = \begin{bmatrix}
        p_k^0 & \zero   & \zero   \\
        \zero & p_k^0   & \zero   \\
        \zero & \zero   & p_k^0
      \end{bmatrix} \in \R^{3m_k\times 3}$;
    \item $\ell=2$ (full affine, $d_2=9$):
      $P_{k,2} = \begin{bmatrix}
        p_k^0 & \zero & \zero & p_k^s & \zero & \zero & p_k^t & \zero & \zero \\[2pt]
        \zero & p_k^0 & \zero & \zero & p_k^s & \zero & \zero & p_k^t & \zero \\[2pt]
        \zero & \zero & p_k^0 & \zero & \zero & p_k^s & \zero & \zero & p_k^t
      \end{bmatrix} \in \R^{3m_k\times 9}$.
  \end{itemize}
  Each column occupies only the block-row matching its direction
  $\beta\in\{n,t_1,t_2\}$.
\end{definition}

The global restriction at level $\ell$ is the block-diagonal assembly
\begin{equation}\label{eq:Pcal}
  \Pcal_\ell := \bigoplus_{k=1}^m P_{k,\ell}
  \;:\; \R^{d_\ell m} \;\longrightarrow\; \Vspace.
\end{equation}
The corresponding \textbf{trace subspace} is
$\Wspace_\ell := \Range(\Pcal_\ell) \subset \Vspace$.

\begin{lemma}[Nested trace subspaces]\label{lem:nested}
  $\Wspace_0 \subseteq \Wspace_1 \subseteq \Wspace_2 \subseteq \Vspace$.
\end{lemma}

\begin{proof}
  Each $P_{k,0}$ (resp.\ $P_{k,1}$) is obtained from $P_{k,1}$
  (resp.\ $P_{k,2}$) by selecting a subset of its columns.
  The block-diagonal structure of $\Pcal_\ell$ preserves the inclusion
  of ranges.
\end{proof}

\begin{remark}[Coefficient-space nesting is not literal]\label{rem:coeff-nest}
  The coefficient spaces $\R^{d_0 m}$, $\R^{d_1 m}$, and $\R^{d_2 m}$
  have different dimensions and are not literally nested unless
  explicit embeddings are specified.  The intrinsic nesting used in
  the Galerkin argument is the nesting of the lifted trace ranges
  $\Wspace_\ell = \Range(\Pcal_\ell) \subseteq \Vspace$.
  All results below are stated in terms of these trace subspaces.
  When we refer to ``level $\ell$'' we mean the subspace $\Wspace_\ell$
  and its associated restriction operator $\Pcal_\ell$.
\end{remark}

% ======================================================================
\section{Best approximation in the energy seminorm}
% ======================================================================

For a fixed level $\ell$, the \textbf{reduced Schur complement} and
\textbf{reduced right-hand side} are
\begin{equation}\label{eq:S-ell}
  \Schur_\ell = \Pcal_\ell^\top \Shatuu\, \Pcal_\ell
  \;\in\; \R^{d_\ell m \times d_\ell m},\qquad
  \bm{F}_\ell = \Pcal_\ell^\top \Fhat.
\end{equation}

\begin{lemma}[Reduced consistency]\label{lem:reduced-cons}
  Under Assumption~\ref{ass:consistency}, for
  every level $\ell$,
  $\bm{F}_\ell \in \Range(\Schur_\ell)$.
\end{lemma}

\begin{proof}
  Let $\bm{\xi}$ be any solution of $\Shatuu\,\bm{\xi} = \Fhat$.
  Set $A_\ell = (\Shatuu)^{1/2}\Pcal_\ell$; then
  $\Schur_\ell = A_\ell^\top A_\ell$ and
  $\bm{F}_\ell = A_\ell^\top (\Shatuu)^{1/2}\bm{\xi}$.
  Finite-dimensional linear algebra gives
  $\Range(A_\ell^\top A_\ell) = \Range(A_\ell^\top)$,
  hence $\bm{F}_\ell \in \Range(\Schur_\ell)$.
\end{proof}

Consequently there exists at least one solution $\trac_\ell \in
\R^{d_\ell m}$ of the reduced system
$\Schur_\ell \trac_\ell = \bm{F}_\ell$.
Define the corresponding C-DD trace approximation
\begin{equation}\label{eq:xi-ell}
  \bm{\xi}_\ell := \Pcal_\ell \trac_\ell \;\in\; \Wspace_\ell \subset \Vspace.
\end{equation}
Let $\bm{\xi} \in \Vspace$ be any fixed solution of the full C-DD
system $\Shatuu\,\bm{\xi} = \Fhat$.

\begin{theorem}[Galerkin best approximation]\label{thm:best}
  Under Assumptions~\ref{ass:local}, \ref{ass:matching}, and
  \ref{ass:consistency}, for each level $\ell\in\{0,1,2\}$,
  \begin{equation}\label{eq:best}
    \|\bm{\xi} - \bm{\xi}_\ell\|_{\Shatuu}
    = \min_{\bm{\mu}\in\R^{d_\ell m}}
      \|\bm{\xi} - \Pcal_\ell\bm{\mu}\|_{\Shatuu}.
  \end{equation}
  When $\Shatuu$ is semidefinite the minimiser need not be unique,
  but the seminorm value is independent of the choice of
  representatives for both $\bm{\xi}$ and $\trac_\ell$.
\end{theorem}

\begin{proof}
  From $\Shatuu\,\bm{\xi} = \Fhat$,
  $\Schur_\ell \trac_\ell = \bm{F}_\ell$,
  $\Schur_\ell = \Pcal_\ell^\top \Shatuu\,\Pcal_\ell$, and
  $\bm{F}_\ell = \Pcal_\ell^\top \Fhat$, we obtain
  \begin{equation}\label{eq:ortho}
    \Pcal_\ell^\top \Shatuu\,(\bm{\xi}_\ell - \bm{\xi}) = \zero.
  \end{equation}
  Hence $\bm{e}_\ell := \bm{\xi} - \bm{\xi}_\ell$ is
  $\Shatuu$-orthogonal to $\Wspace_\ell = \Range(\Pcal_\ell)$.

  For arbitrary $\bm{\mu}\in\R^{d_\ell m}$, decompose
  $\bm{\xi} - \Pcal_\ell\bm{\mu}
   = \bm{e}_\ell + \bm{v}$
  with $\bm{v} := \bm{\xi}_\ell - \Pcal_\ell\bm{\mu}
  \in \Wspace_\ell$.
  By~\eqref{eq:ortho}, $\bm{e}_\ell^\top \Shatuu\,\bm{v} = 0$, and the
  Pythagorean identity in the $\Shatuu$ seminorm yields
  $\|\bm{\xi} - \Pcal_\ell\bm{\mu}\|_{\Shatuu}^2
   = \|\bm{e}_\ell\|_{\Shatuu}^2 + \|\bm{v}\|_{\Shatuu}^2
   \ge \|\bm{e}_\ell\|_{\Shatuu}^2$.
  Equality at $\bm{\mu} = \trac_\ell$ proves~\eqref{eq:best}.

  Semidefinite independence: if $\trac_\ell,\trac_\ell'$ both solve
  $\Schur_\ell\trac = \bm{F}_\ell$, then
  $\Pcal_\ell(\trac_\ell-\trac_\ell') \in \Ker(\Shatuu)$,
  so the seminorm on both sides of~\eqref{eq:best} is unchanged.
  Similarly, any two full solutions differ by an element of
  $\Ker(\Shatuu)$, and adding an element of this kernel does not change
  any $\Shatuu$-seminorm distance.
\end{proof}

% ======================================================================
\section{Nested orthogonal decomposition}
% ======================================================================

\begin{theorem}[Orthogonal decomposition for nested trace subspaces]\label{thm:nested}
  Under the same assumptions as Theorem~\ref{thm:best}, let
  $\ell < \ell'$, so that $\Wspace_\ell \subseteq \Wspace_{\ell'}$ by
  Lemma~\ref{lem:nested}.  Then
  \begin{equation}\label{eq:nested}
    \|\bm{\xi} - \bm{\xi}_\ell\|_{\Shatuu}^2
    = \|\bm{\xi} - \bm{\xi}_{\ell'}\|_{\Shatuu}^2
      + \|\bm{\xi}_{\ell'} - \bm{\xi}_\ell\|_{\Shatuu}^2.
  \end{equation}
\end{theorem}

\begin{proof}
  By Theorem~\ref{thm:best} at level $\ell'$,
  $\bm{e}_{\ell'} = \bm{\xi} - \bm{\xi}_{\ell'}$ is
  $\Shatuu$-orthogonal to $\Wspace_{\ell'}$.
  The difference $\bm{\xi}_{\ell'} - \bm{\xi}_\ell$ satisfies
  $\bm{\xi}_{\ell'} \in \Wspace_{\ell'}$ and
  $\bm{\xi}_\ell \in \Wspace_\ell \subseteq \Wspace_{\ell'}$,
  hence $\bm{\xi}_{\ell'} - \bm{\xi}_\ell \in \Wspace_{\ell'}$.
  Therefore $\bm{e}_{\ell'} \perp_{\Shatuu}
  (\bm{\xi}_{\ell'} - \bm{\xi}_\ell)$.
  The decomposition
  $\bm{\xi} - \bm{\xi}_\ell = \bm{e}_{\ell'} +
   (\bm{\xi}_{\ell'} - \bm{\xi}_\ell)$
  is $\Shatuu$-orthogonal; the Pythagorean identity gives~\eqref{eq:nested}.
\end{proof}

\begin{corollary}[Three-level gain decomposition]\label{cor:gains}
  For $\Wspace_0 \subseteq \Wspace_1 \subseteq \Wspace_2$,
  \begin{equation}\label{eq:gains}
    \|\bm{\xi} - \bm{\xi}_0\|_{\Shatuu}^2
    = \|\bm{\xi} - \bm{\xi}_2\|_{\Shatuu}^2
      + \|\bm{\xi}_2 - \bm{\xi}_1\|_{\Shatuu}^2
      + \|\bm{\xi}_1 - \bm{\xi}_0\|_{\Shatuu}^2.
  \end{equation}
  Interpretation of the three right-hand-side terms:
  \begin{itemize}
    \item $\|\bm{\xi} - \bm{\xi}_2\|_{\Shatuu}^2$:
      \textbf{affine residual} --- error after the full nine-mode
      affine approximation;
    \item $\|\bm{\xi}_2 - \bm{\xi}_1\|_{\Shatuu}^2$:
      \textbf{linear gain} --- reduction from adding the six linear
      modes $\{s,t\}\times\{n,t_1,t_2\}$ to the constant-vector space;
    \item $\|\bm{\xi}_1 - \bm{\xi}_0\|_{\Shatuu}^2$:
      \textbf{tangential gain} --- reduction from adding the two
      tangential constant modes to the single normal constant mode.
  \end{itemize}
\end{corollary}

\begin{proof}
  Apply Theorem~\ref{thm:nested} with $(\ell,\ell')=(0,1)$ and
  $(\ell,\ell')=(1,2)$, then substitute.
\end{proof}

\begin{remark}[Four-space diagnostic diamond]\label{rem:four-space}
  The three-level chain $\Wspace_0\subseteq\Wspace_1\subseteq\Wspace_2$
  combines all six linear modes into one gain term.  A controlled
  diagnostic can instead use four global trace subspaces whose
  per-interface mode contents are
  \begin{align*}
    \Wspace_{0n} &: \{1\cdot n\},\\
    \Wspace_{0v} &: \{1\cdot n,\;1\cdot t_1,\;1\cdot t_2\},\\
    \Wspace_{1n} &: \{1\cdot n,\;s\cdot n,\;t\cdot n\},\\
    \Wspace_{1v} &: \{1,s,t\}\times\{n,t_1,t_2\}.
  \end{align*}
  Here each notation denotes the range of the corresponding
  block-diagonal modal lifting assembled over all interfaces, and
  $\bm{\xi}_{0n}$, $\bm{\xi}_{0v}$, $\bm{\xi}_{1n}$, and
  $\bm{\xi}_{1v}$ denote the associated Galerkin trace solutions.
  The following inclusions always hold:
  $\Wspace_{0n}\subseteq\Wspace_{0v}\subseteq\Wspace_{1v}$
  and
  $\Wspace_{0n}\subseteq\Wspace_{1n}\subseteq\Wspace_{1v}$.
  Whether $\Wspace_{0v}$ and $\Wspace_{1n}$ are nonnested depends on
  the discrete trace space.  They are nonnested, for instance, if there
  exists at least one interface on which the projected linear shapes
  $\{p_k^s,p_k^t\}$ contain a direction outside
  $\operatorname{span}\{p_k^0\}$ and the constant tangential liftings are
  nonzero.  The following two nested-path Pythagorean decompositions are
  valid irrespective of whether either inclusion is strict.  Applying
  Theorem~\ref{thm:nested} gives
  \[
  \begin{aligned}
  \|\bm{\xi}-\bm{\xi}_{0n}\|_{\Shatuu}^2
  &= \|\bm{\xi}-\bm{\xi}_{1v}\|_{\Shatuu}^2
   + \|\bm{\xi}_{1v}-\bm{\xi}_{0v}\|_{\Shatuu}^2
   + \|\bm{\xi}_{0v}-\bm{\xi}_{0n}\|_{\Shatuu}^2,\\
  &= \|\bm{\xi}-\bm{\xi}_{1v}\|_{\Shatuu}^2
   + \|\bm{\xi}_{1v}-\bm{\xi}_{1n}\|_{\Shatuu}^2
   + \|\bm{\xi}_{1n}-\bm{\xi}_{0n}\|_{\Shatuu}^2.
  \end{aligned}
  \]
  These are path-dependent marginal-gain decompositions, not a unique
  causal attribution.  Two cleanly defined error-reduction contrasts
  sharing the same baseline $\bm{\xi}_{0n}$ are obtained from the
  Pythagorean identity (Theorem~\ref{thm:nested}):
  \[
    \Delta_T
    := \|\bm{\xi}-\bm{\xi}_{0n}\|_{\Shatuu}^2
       -\|\bm{\xi}-\bm{\xi}_{0v}\|_{\Shatuu}^2
    = \|\bm{\xi}_{0v}-\bm{\xi}_{0n}\|_{\Shatuu}^2 \ge 0,
  \]
  \[
    \Delta_N
    := \|\bm{\xi}-\bm{\xi}_{0n}\|_{\Shatuu}^2
       -\|\bm{\xi}-\bm{\xi}_{1n}\|_{\Shatuu}^2
    = \|\bm{\xi}_{1n}-\bm{\xi}_{0n}\|_{\Shatuu}^2 \ge 0.
  \]
  These quantify the reductions in squared Schur-energy error under the
  two specified nested enrichments.  The complementary path-conditional
  gains are
  \[
    \Delta_{N\mid T}
    := \|\bm{\xi}_{1v}-\bm{\xi}_{0v}\|_{\Shatuu}^2,\qquad
    \Delta_{T\mid N}
    := \|\bm{\xi}_{1v}-\bm{\xi}_{1n}\|_{\Shatuu}^2.
  \]
  The total reduction from the classic baseline to the full affine
  solution therefore satisfies
  \[
    \|\bm{\xi}-\bm{\xi}_{0n}\|_{\Shatuu}^2
    -\|\bm{\xi}-\bm{\xi}_{1v}\|_{\Shatuu}^2
    = \Delta_T + \Delta_{N\mid T}
    = \Delta_N + \Delta_{T\mid N}.
  \]
  In general $\Delta_N \neq \Delta_{N\mid T}$ and
  $\Delta_T \neq \Delta_{T\mid N}$, illustrating that
  the two paths produce path-dependent marginal-gain decompositions.
  No identity of the form
  $\Delta_{\mathrm{total}} = \Delta_T + \Delta_N$ is asserted, and no
  per-interface additive causal decomposition is claimed.
\end{remark}

% ======================================================================
\section{Velocity reconstruction and trace error identity}
% ======================================================================

\subsection{Velocity reconstruction from C-DD trace data}

For a C-DD trace vector $\bm{\zeta}\in\Vspace$ (unknown internal DOFs)
and prescribed boundary traction data $\bm{\xi}^k$, the local velocity
field on pore $i$ is
\begin{equation}\label{eq:reconstruct}
  U_i(\bm{\zeta};\,\bm{\xi}^k)
  = -H_i\Bigl[(N_i^u)^\top \widehat{R}_i^u \bm{\zeta}
               + (N_i^k)^\top \widehat{R}_i^k \bm{\xi}^k\Bigr]
  \;\in\; \R^{n_{u,i}}.
\end{equation}
The first term couples the unknown interface DOFs to the pore
interior; the second term injects the prescribed boundary tractions.
The global reconstructed velocity field $\vel_h(\bm{\zeta})$ is the
discontinuous assembly of the local fields $\{U_i(\bm{\zeta};\,
\bm{\xi}^k)\}$ across pores.

\begin{lemma}[Trace-to-velocity error identity]\label{lem:vel-id}
  For any two trace vectors $\bm{\zeta},\bm{\zeta}'\in\Vspace$ sharing
  the same prescribed boundary data $\bm{\xi}^k$,
  \begin{equation}\label{eq:vel-id}
    \sum_{i=1}^N \|U_i(\bm{\zeta}) - U_i(\bm{\zeta}')\|_{A_i}^2
    = \|\bm{\zeta} - \bm{\zeta}'\|_{\Shatuu}^2,
  \end{equation}
  where $\|U\|_{A_i}^2 := U^\top A_i U$.
\end{lemma}

\begin{proof}
  The boundary coupling terms $(N_i^k)^\top \widehat{R}_i^k \bm{\xi}^k$
  are identical for both arguments and cancel in the difference:
  \begin{align*}
    U_i(\bm{\zeta}) - U_i(\bm{\zeta}')
    &= -H_i (N_i^u)^\top \widehat{R}_i^u (\bm{\zeta} - \bm{\zeta}').
  \end{align*}
  Using $H_i^\top A_i H_i = H_i$,
  \begin{align*}
    \|U_i(\bm{\zeta}) - U_i(\bm{\zeta}')\|_{A_i}^2
    &= (\bm{\zeta}-\bm{\zeta}')^\top \widehat{R}_i^{u\top}
       N_i^u H_i A_i H_i (N_i^u)^\top \widehat{R}_i^u
       (\bm{\zeta}-\bm{\zeta}') \\
    &= (\bm{\zeta}-\bm{\zeta}')^\top \widehat{R}_i^{u\top}
       \Shatuu_i \widehat{R}_i^u (\bm{\zeta}-\bm{\zeta}').
  \end{align*}
  Summing over $i$ and using~\eqref{eq:Shat} gives~\eqref{eq:vel-id}.
\end{proof}

\subsection{Notation for the true solution and its approximations}

Let $\vel_h := \vel_h(\bm{\xi})$ be the C-DD full-trace reconstruction
using a solution $\bm{\xi}$ of~\eqref{eq:full-sys}.
For level $\ell$, let $\vel_\ell := \vel_h(\bm{\xi}_\ell)$.
Let $\vel$ denote the exact Stokes solution with the prescribed
continuous boundary-traction data.  The full C-DD reconstruction
$\vel_h$ and the reduced reconstruction $\vel_\ell$ share the same
prescribed discrete boundary-traction vector $\bm{\xi}^k$, obtained
from the boundary restriction in~\eqref{eq:Fhat}.  Consequently the
discrete boundary terms cancel exactly in $\vel_h-\vel_\ell$.
Any discrepancy between the continuous boundary data and its discrete
representation is included in the C-DD discretisation error
$\vel-\vel_h$.
All velocity error fields considered below satisfy homogeneous
wall data.

Let $\Gamma_{w,i}$ denote the no-slip wall portion of
$\partial\pore{i}$ and define the broken wall-zero space
\[
  X_0:=\prod_{i=1}^N
  \{\bm{v}_i\in H^1(\pore{i})^3:\ \bm{v}_i|_{\Gamma_{w,i}}=\zero\}.
\]
All velocity error fields used below belong to $X_0$ because the exact,
full C-DD, and reduced solutions impose the same no-slip wall data.
For a (possibly discontinuous) field $\bm{v}$ define
\begin{align}
  \|\bm{v}\|_A^2 &:= \sum_{i=1}^N \nu\int_{\pore{i}}
                     \nabla\bm{v}:\nabla\bm{v}\,\mathrm{d}x,
                     \label{eq:A-norm} \\[4pt]
  \|\bm{v}\|_{L^2}^2 &:= \sum_{i=1}^N \int_{\pore{i}}
                         |\bm{v}|^2\,\mathrm{d}x, \label{eq:L2-norm} \\[4pt]
  |\bm{v}|_{H^1}^2 &:= \sum_{i=1}^N \int_{\pore{i}}
                        \nabla\bm{v}:\nabla\bm{v}\,\mathrm{d}x
                        = \frac{1}{\nu}\|\bm{v}\|_A^2. \label{eq:H1-semi}
\end{align}
On the full broken product space, the two gradient-based quantities
$\|\cdot\|_A$ and $|\cdot|_{H^1}$ are seminorms.  If every $\pore{i}$ is connected and Lipschitz and
$|\Gamma_{w,i}|>0$, Poincar\'e--Friedrichs makes $\|\cdot\|_A$ a norm
on $X_0$.  These conditions are assumed whenever the term ``$A$-norm''
is used below.  For a discrete function $\bm{v}_h$ with local coefficient
vectors $\{V_i\}$,
$\|\bm{v}_h\|_A^2 = \sum_i V_i^\top A_i V_i$.

% ======================================================================
\section{Volume error---trace error separation}
% ======================================================================

\subsection{Two-sided $A$-norm separation and the trace-dominated regime}

Set
\begin{equation}\label{eq:DE}
  D_h := \|\vel - \vel_h\|_A,\qquad
  E_\ell := \|\vel_h - \vel_\ell\|_A
          = \|\bm{\xi} - \bm{\xi}_\ell\|_{\Shatuu},
\end{equation}
where the second equality is Lemma~\ref{lem:vel-id} with
$\bm{\zeta}=\bm{\xi}$, $\bm{\zeta}'=\bm{\xi}_\ell$.

\begin{theorem}[Two-sided volume-trace bound: $A$-norm]\label{thm:sep-A}
  Under Assumptions~\ref{ass:local}, \ref{ass:matching}, and
  \ref{ass:consistency},
  for each level $\ell$,
  \begin{equation}\label{eq:sep-A}
    |E_\ell - D_h|
    \;\le\; \|\vel - \vel_\ell\|_A
    \;\le\; E_\ell + D_h.
  \end{equation}
  Equivalently,
  $|\,\|\vel-\vel_\ell\|_A - E_\ell\,| \le D_h$.
  Hence, whenever $E_\ell>0$,
  \begin{equation}\label{eq:relative}
    \Bigl|
      \frac{\|\vel-\vel_\ell\|_A}{E_\ell} - 1
    \Bigr|
    \le \frac{D_h}{E_\ell}.
  \end{equation}
\end{theorem}

\begin{proof}
  By the reverse and forward triangle inequalities for the $A$-norm,
  \[
    \bigl|\|\vel_h - \vel_\ell\|_A - \|\vel - \vel_h\|_A\bigr|
    \le \|\vel - \vel_\ell\|_A
    \le \|\vel_h - \vel_\ell\|_A + \|\vel - \vel_h\|_A.
  \]
  Substituting $E_\ell$ and $D_h$ from~\eqref{eq:DE} gives the first
  statement; dividing by $E_\ell>0$ yields~\eqref{eq:relative}.
\end{proof}

\begin{remark}[Fixed-mesh trace domination and its limitation]\label{rem:floor}
  On a fixed mesh $h$, define the discretisation-to-reduction ratio
  \begin{equation}\label{eq:rho}
    \rho_{\ell,h} := \frac{D_h}{E_{\ell,h}},
  \end{equation}
  which is well-defined whenever $E_{\ell,h}>0$.
  For a prescribed tolerance $0\le\eta<1$, we call the discretisation
  \textbf{trace-dominated at level $\eta$} if $\rho_{\ell,h}\le\eta$.
  In that case~\eqref{eq:relative} gives the two-sided estimate
  \begin{equation}\label{eq:trace-dom}
    (1-\eta)E_{\ell,h}
    \;\le\; \|\vel - \vel_{\ell,h}\|_A
    \;\le\; (1+\eta)E_{\ell,h}.
  \end{equation}

  This statement is strictly mesh-dependent.  It does not imply that
  mesh refinement leaves the total error unchanged, because both
  $D_h$ and $E_{\ell,h}$ vary with $h$.  A positive asymptotic lower
  bound for the total $A$-error requires control of the
  discretisation error as well.  From~\eqref{eq:sep-A},
  $\|\vel-\vel_{\ell,h}\|_A \ge E_{\ell,h}-D_h$, hence
  \[
    \liminf_{h\to0}\,\|\vel-\vel_{\ell,h}\|_A
    \;\ge\; \liminf_{h\to0}E_{\ell,h} - \limsup_{h\to0}D_h.
  \]
  A sufficient condition for a positive asymptotic error floor is
  therefore
  \begin{equation}\label{eq:floor-cond}
    \liminf_{h\to0}E_{\ell,h} \;>\; \limsup_{h\to0}D_h.
  \end{equation}
  In particular, $D_h\to0$ together with
  $\liminf_{h\to0}E_{\ell,h}=c_*>0$ is sufficient.
  The present theory does not establish whether~\eqref{eq:floor-cond}
  holds for the geometries under study; numerical plateaus may motivate
  such a hypothesis but are not a consequence of these theorems.
  Moreover, the planar-interface theory developed here does not justify
  interpreting a plateau on a non-planar watershed partition as
  evidence for~\eqref{eq:floor-cond}.

  \textbf{Computability.}  $E_{\ell,h}
  =\|\bm{\xi}-\bm{\xi}_\ell\|_{\Shatuu}$ is computable when the full
  C-DD and reduced discrete solutions are both available.  In contrast,
  $D_h=\|\vel-\vel_h\|_A$ requires the exact PDE solution and is not
  directly computable in a generic problem; it must be replaced by a
  manufactured-solution error, an independent overkill reference, or a
  separate discretisation-error estimator.
\end{remark}

\subsection{Analogous bounds in $L^2$ and broken $H^1$}

\begin{theorem}[Two-sided bound: $L^2$-norm]\label{thm:sep-L2}
  Under the same assumptions, suppose each pore $\pore{i}$ admits a
  Poincar\'e--Friedrichs constant $C_{P,i}>0$ such that
  $\|\bm{v}\|_{L^2(\pore{i})} \le C_{P,i}|\bm{v}|_{H^1(\pore{i})}$
  for all $\bm{v}$ vanishing on the wall portion of
  $\partial\pore{i}$.  Set $C_P := \max_i C_{P,i}$ and define
  \begin{equation}\label{eq:TL2}
    T_\ell^{L^2} := \|\vel_h - \vel_\ell\|_{L^2},\qquad
    D_h^{L^2} := \|\vel - \vel_h\|_{L^2}.
  \end{equation}
  Then the following hold.
  \begin{enumerate}[label=(\alph*),nosep,leftmargin=*]
    \item \textbf{Two-sided bound with $T_\ell^{L^2}$:}
      $|T_\ell^{L^2} - D_h^{L^2}|
       \le \|\vel - \vel_\ell\|_{L^2}
       \le T_\ell^{L^2} + D_h^{L^2}$.
    \item \textbf{Poincar\'e upper bound:}
      $T_\ell^{L^2} \le \frac{C_P}{\sqrt{\nu}}\,E_\ell$, and therefore
      $\|\vel - \vel_\ell\|_{L^2}
       \le D_h^{L^2} + \frac{C_P}{\sqrt{\nu}}\,E_\ell$.
  \end{enumerate}
  No corresponding lower bound in terms of $E_\ell$ is available without
  a reverse stability estimate (e.g.\ an inverse inequality bounding
  the $A$-norm from above by the $L^2$-norm on the discrete velocity
  space).
\end{theorem}

\begin{proof}
  (a)~Apply the reverse and forward triangle inequalities for the
  $L^2$-norm to $\vel-\vel_\ell = (\vel-\vel_h) + (\vel_h-\vel_\ell)$.
  (b)~From Lemma~\ref{lem:vel-id} and per-pore Poincar\'e--Friedrichs,
  $(T_\ell^{L^2})^2
   = \sum_i \|U_i(\bm{\xi})-U_i(\bm{\xi}_\ell)\|_{L^2(\pore{i})}^2
   \le \frac{C_P^2}{\nu}\sum_i \|U_i(\bm{\xi})-U_i(\bm{\xi}_\ell)\|_{A_i}^2
   = \frac{C_P^2}{\nu}\,E_\ell^2$.
  Taking square roots and substituting into the upper bound of (a)
  gives the second statement.
\end{proof}

\begin{theorem}[Two-sided bound: broken-$H^1$ seminorm]\label{thm:sep-H1}
  Under the same assumptions, define
  $T_\ell^{H^1} := |\vel_h - \vel_\ell|_{H^1}$ and
  $D_h^{H^1} := |\vel - \vel_h|_{H^1}$.
  From~\eqref{eq:H1-semi} and Lemma~\ref{lem:vel-id},
  $T_\ell^{H^1} = E_\ell/\sqrt{\nu}$ exactly.
  Then
  \begin{equation}\label{eq:sep-H1}
    |T_\ell^{H^1} - D_h^{H^1}|
    \;\le\; |\vel - \vel_\ell|_{H^1}
    \;\le\; T_\ell^{H^1} + D_h^{H^1}
    \;=\; \frac{1}{\sqrt{\nu}}\,E_\ell + D_h^{H^1}.
  \end{equation}
\end{theorem}

\begin{proof}
  The identity $T_\ell^{H^1} = E_\ell/\sqrt{\nu}$ follows from
  $|\bm{v}|_{H^1}^2 = \|\bm{v}\|_A^2/\nu$ and Lemma~\ref{lem:vel-id}.
  The seminorm triangle inequalities give the two-sided bound.
  Unlike the $L^2$ case, both the upper and the lower bound are
  available here because $T_\ell^{H^1}$ is exactly computable
  from $E_\ell$ (no additional inequality needed).
\end{proof}

\begin{remark}[Pessimism of the $L^2$ constant]\label{rem:CP}
  The constant $C_P=\max_i C_{P,i}$ is a worst-case Poincar\'e constant
  and may be pessimistic.  For the particular reduction error
  $\bm{e}_\ell:=\vel_h-\vel_\ell$, the field-dependent effective
  Poincar\'e ratio is
  \[
    C_{P,\mathrm{eff}}(\bm{e}_\ell)
      :=\frac{\|\bm{e}_\ell\|_{L^2}}{|\bm{e}_\ell|_{H^1}}
      =\sqrt{\nu}\,\frac{\|\bm{e}_\ell\|_{L^2}}
                              {\|\bm{e}_\ell\|_A},
  \]
  whenever $\bm{e}_\ell\neq0$.  Its square is a generalized Rayleigh
  quotient for the mass and stiffness forms evaluated at this specific
  error field; it is not a uniform stability constant for the whole
  space.
\end{remark}

% ======================================================================
\section{Discussion}
% ======================================================================

\subsection{Summary of the theoretical chain}

\begin{enumerate}[leftmargin=*,itemsep=4pt]
  \item \textbf{Best approximation} (Theorem~\ref{thm:best}).
    Under the explicit consistency assumption for the full C-DD Schur
    system, the level-$\ell$ trace solution is a
    $\Shatuu$-seminorm best approximation to the full C-DD trace in
    $\Wspace_\ell$.
  \item \textbf{Nested decomposition} (Theorem~\ref{thm:nested},
    Corollary~\ref{cor:gains}).
    For nested modal trace spaces, squared Schur-energy errors satisfy
    exact Pythagorean identities.  The three-level chain therefore
    separates the classic-P0 reduction error into the residual after
    full affine approximation and the gains associated with the two
    nested enrichments.
  \item \textbf{Volume--trace separation}.
    In the broken $A$ measure,
    \[
      |E_\ell-D_h|\le\|\vel-\vel_\ell\|_A\le E_\ell+D_h.
    \]
    In the broken-$H^1$ seminorm the same statement holds after the
    exact scaling $E_\ell/\sqrt{\nu}$.  In $L^2$, the exact two-sided
    triangle bound involves the actual reduction error
    $T_\ell^{L^2}=\|\vel_h-\vel_\ell\|_{L^2}$; the Schur-energy error
    $E_\ell$ supplies only the upper estimate
    $T_\ell^{L^2}\le(C_P/\sqrt{\nu})E_\ell$ unless an additional
    reverse stability estimate is available.
\end{enumerate}

\subsection{Limitations of the current theory}

\begin{enumerate}[leftmargin=*,itemsep=4pt]
  \item \textbf{Full C-DD consistency is assumed, not re-proved here.}
    Assumption~\ref{ass:consistency} is sufficient for all Galerkin
    projection arguments in this document.  A self-contained derivation
    would start from the complete broken-pressure mixed Stokes system
    and prove that its elimination produces a consistent Schur system.

  \item \textbf{No C-DD--monolithic equivalence theorem is claimed.}
    The C-DD reference uses pressure that is discontinuous across pore
    interfaces, whereas the usual monolithic Taylor--Hood discretisation
    uses globally continuous $P_1$ pressure.  Equality of their discrete
    velocity solutions is not established.  Even if such an equality
    were proved, a separate finite-element discretisation error between
    the monolithic discrete solution and the exact PDE solution would
    remain.

  \item \textbf{No mesh-independent positive error floor is proved.}
    Theorem~\ref{thm:sep-A} is a fixed-mesh separation result.  To infer
    a refinement-resistant positive floor one would additionally need
    to control the $h$-dependence of $E_{\ell,h}$ and $D_h$; a sufficient
    condition is $\liminf_{h\to0}E_{\ell,h} > \limsup_{h\to0}D_h$
    (see~\eqref{eq:floor-cond}).  The present theory does not establish
    whether this condition holds.

  \item \textbf{The $L^2$ trace-energy estimate is one-sided.}
    The Poincar\'e--Friedrichs inequality bounds
    $\|\vel_h-\vel_\ell\|_{L^2}$ from above by $E_\ell$, but it yields no
    lower bound in terms of $E_\ell$.  Consequently $E_\ell$ should not
    be called an $L^2$ error floor without an additional reverse
    stability estimate.

  \item \textbf{The affine residual need not be strictly positive.}
    Although the nine-mode affine trace space is generally lower
    dimensional than a nodal C-DD trace space, strict inclusion alone
    does not imply $E_2>0$ for a particular solution.  A positive
    affine residual requires
    $\bm{\xi}\notin\Wspace_2+\Ker(\Shatuu)$.

  \item \textbf{The four-space diamond is diagnostic, not a unique
    causal decomposition.}
    When $\Wspace_{0v}$ and $\Wspace_{1n}$ are nonnested, as under the
    sufficient nondegeneracy condition stated in
    Remark~\ref{rem:four-space}, the two Pythagorean paths provide
    path-indexed marginal-gain decompositions, which are path-dependent
    in general.  Regardless of whether the two middle spaces happen to
    be comparable for a particular discrete trace space, the quantities
    $\Delta_T$ and $\Delta_N$ are controlled error-reduction contrasts
    associated with specified enrichments, not a unique causal
    attribution.

  \item \textbf{Non-planar interfaces are outside the present scope.}
    Every theorem uses a fixed orthonormal frame on a planar interface.
    Faceted or genuinely curved interfaces with spatially varying
    normals introduce a geometric representation issue not covered by
    these results.
\end{enumerate}

% ======================================================================
\begin{thebibliography}{9}

\bibitem{sun2025ddpnm}
S.~Sun, Z.~Wang, L.~Zhang, J.~Zhao.
\newblock A domain decomposition approach to pore-network modeling of porous
  media flow.
\newblock \textit{arXiv:2510.13429v2}, 2026.

\bibitem{toselli2005}
A.~Toselli and O.~Widlund.
\newblock \textit{Domain Decomposition Methods---Algorithms and Theory}.
\newblock Springer, 2005.

\bibitem{horn}
R.~A.~Horn and C.~R.~Johnson.
\newblock \textit{Matrix Analysis}, 2nd ed.
\newblock Cambridge University Press, 2012.

\end{thebibliography}

\end{document}
